\documentclass[10pt,a4paper]{article}

\usepackage[utf8]{inputenc}
\usepackage[T1]{fontenc}
\usepackage[margin=1.5cm]{geometry}
\usepackage{amsmath,amssymb}
\usepackage{xcolor}
\usepackage{tcolorbox}
\tcbuselibrary{breakable}
\usepackage{listings}

\lstset{
    language=C,
    basicstyle=\ttfamily\scriptsize,
    backgroundcolor=\color{gray!10},
    keywordstyle=\color{blue},
    commentstyle=\color{gray},
    stringstyle=\color{red},
    numbers=left,
    numberstyle=\tiny,
    breaklines=true,
    frame=single,
    showstringspaces=false,
    tabsize=4,
    inputencoding=utf8,
    extendedchars=true,
    literate={é}{{\'e}}1 {è}{{\`e}}1 {à}{{\`a}}1 {ç}{{\c{c}}}1 
             {ê}{{\^e}}1 {î}{{\^i}}1 {ô}{{\^o}}1 {û}{{\^u}}1 {ù}{{\`u}}1
             {É}{{\'E}}1 {È}{{\`E}}1 {À}{{\`A}}1 {Ç}{{\c{C}}}1
}

\newtcolorbox{exercicebox}[1]{
    colback=green!5,colframe=green!50!black,
    fonttitle=\bfseries,title=#1,breakable
}

\newtcolorbox{warningbox}[1]{
    colback=red!5,colframe=red!75!black,
    fonttitle=\bfseries,title=#1,breakable
}

\title{\textbf{TD Chapitre 0 : Fondamentaux du C}\\
\large Analyse et Prédiction}
\author{Dr. El Hadji Bassirou TOURÉ\\
Département de Mathématiques et Informatique\\
Université Cheikh Anta Diop de Dakar}
\date{}

\begin{document}

\maketitle

\section*{Exercice 1 : Analyse Multi-Concepts}

\begin{exercicebox}{Types, conversions, opérateurs et évaluation}
Pour chaque programme, prédire la sortie EXACTE et justifier en mobilisant les concepts pertinents (types, conversions, précédence, court-circuit, UB).

\textbf{Programme A:}
\begin{lstlisting}
#include <stdio.h>
int main() {
    char c = 200;
    unsigned char uc = 200;
    int x = c + uc;
    printf("%d %d %d\n", c, uc, x);
    
    int si = -1;
    unsigned int ui = 1;
    printf("%d\n", si < ui);
    printf("%u\n", si);
    return 0;
}
\end{lstlisting}

\textbf{Programme B:}
\begin{lstlisting}
#include <stdio.h>
int main() {
    int a = 5, b = 10, c = 15;
    int r1 = a + b * c;
    int r2 = a & b == 10;
    int r3 = 1 << 2 + 3;
    printf("%d %d %d\n", r1, r2, r3);
    
    unsigned char x = 0xA5;
    printf("%02X %02X %02X\n", x & 0x0F, x >> 4, x ^ 0xFF);
    return 0;
}
\end{lstlisting}

\textbf{Programme C:}
\begin{lstlisting}
#include <stdio.h>
int main() {
    int a = 0, b = 5, c = 0;
    
    if (a++ || b++ && c++) 
        printf("vrai ");
    printf("a=%d b=%d c=%d\n", a, b, c);
    
    int x = 10, y = 0;
    int z = (y != 0) && (x / y > 2);
    printf("z=%d\n", z);
    return 0;
}
\end{lstlisting}
\end{exercicebox}

\section*{Exercice 2 : Variables, Stockage et Linkage}

\begin{exercicebox}{static, extern, portée et durée de vie}
\textbf{A.} Prédire la sortie exacte :
\begin{lstlisting}
#include <stdio.h>
int x = 1;
int f(void) { static int x = 10; return x++; }
int g(void) { int x = 100; return x++; }
int h(void) { return x++; }

int main() {
    printf("%d ", f()); printf("%d ", f()); printf("%d\n", f());
    printf("%d ", g()); printf("%d ", g()); printf("%d\n", g());
    printf("%d ", h()); printf("%d ", h()); printf("%d\n", h());
    printf("x=%d\n", x);
    return 0;
}
\end{lstlisting}

\textbf{B.} Soit deux fichiers compilés avec \texttt{gcc a.c b.c -o prog}. Pour chaque cas, indiquer si ça compile, linke, et si oui quelle est la sortie.

\textbf{Cas 1:}
\begin{lstlisting}
// a.c                          // b.c
#include <stdio.h>              int val = 42;
extern int val;
int main() { printf("%d\n", val); return 0; }
\end{lstlisting}

\textbf{Cas 2:}
\begin{lstlisting}
// a.c                          // b.c
#include <stdio.h>              static int val = 42;
extern int val;
int main() { printf("%d\n", val); return 0; }
\end{lstlisting}

\textbf{Cas 3:}
\begin{lstlisting}
// a.c                          // b.c
#include <stdio.h>              int val = 42;
int val = 100;
int main() { printf("%d\n", val); return 0; }
\end{lstlisting}
\end{exercicebox}

\section*{Exercice 3 : Macros, sizeof et Pièges}

\begin{exercicebox}{Préprocesseur, sizeof, évaluation}
\textbf{A.} Prédire chaque sortie en expandant manuellement les macros :
\begin{lstlisting}
#include <stdio.h>
#define A 2
#define B A + 3
#define C B * 2
#define SQ(x) x * x
#define SAFE_SQ(x) ((x) * (x))

int main() {
    printf("%d %d\n", B, C);
    printf("%d %d\n", 100 / C, 100 / B);
    printf("%d %d\n", SQ(1+2), SAFE_SQ(1+2));
    return 0;
}
\end{lstlisting}

\textbf{B.} Prédire la sortie et expliquer :
\begin{lstlisting}
#include <stdio.h>
int f(void) { printf("!"); return 1; }

int main() {
    int x = 10;
    size_t s1 = sizeof(x = 999);
    size_t s2 = sizeof(x++);
    size_t s3 = sizeof(f());
    printf("x=%d s1=%zu s2=%zu s3=%zu\n", x, s1, s2, s3);
    return 0;
}
\end{lstlisting}

\textbf{C.} Expliquer les résultats (x86-64) :
\begin{lstlisting}
#include <stdio.h>
int main() {
    printf("%zu %zu\n", sizeof('A'), sizeof(3.14));
    return 0;
}
\end{lstlisting}
\end{exercicebox}

\section*{Exercice 4 : Contrôle et Comportement Indéfini}

\begin{exercicebox}{Switch, boucles, goto, UB}
\textbf{A.} Prédire la sortie :
\begin{lstlisting}
#include <stdio.h>
int main() {
    for (int x = 0; x < 3; x++) {
        switch (x) {
            case 0: printf("A");
            case 1: printf("B"); break;
            case 2: printf("C");
            default: printf("D");
        }
        printf(" ");
    }
    return 0;
}
\end{lstlisting}

\textbf{B.} Pour chaque fragment : UB ou défini ? Si défini, quelle valeur ?
\begin{lstlisting}
// (1)                    // (2)                    // (3)
int a = 10;               #include <limits.h>       #include <limits.h>
int b = (a++, ++a, a+1);  int c = INT_MAX + 1;      unsigned u = UINT_MAX + 1;

// (4)                    // (5)
int d = 1 << 31;          int e; printf("%d", e);
\end{lstlisting}

\textbf{C.} Réécrire avec goto (pattern sortie de boucles imbriquées) :
\begin{lstlisting}
// Trouver i,j tels que i*i + j*j == 50
int found = 0, ri = 0, rj = 0;
for (int i = 1; i <= 10 && !found; i++) {
    for (int j = 1; j <= 10 && !found; j++) {
        if (i*i + j*j == 50) { found = 1; ri = i; rj = j; }
    }
}
\end{lstlisting}
\end{exercicebox}

\begin{warningbox}{Consignes}
Pour chaque question : identifier les concepts en jeu, expliquer le raisonnement, donner la sortie exacte ou justifier pourquoi c'est indéfini.
\end{warningbox}

\end{document}